\renewcommand{\Anno}{2022}

\renewcommand{\NumeroProblema}{2}

\renewcommand{\Testo}{
Il Marathon Club Sant'Anna ha 100 iscritti, più il Presidente Antonio. Tutte le settimane, i componenti ed il Presidente si ritrovano per un allenamento di gruppo lungo un percorso della lunghezza di 10~km.

\begin{enumerate}
\renewcommand{\labelenumi}{(\alph{enumi})}

\item  Nel periodo iniziale della pandemia, per evitare assembramenti, i componenti del gruppo hanno deciso di effettuare l'allenamento partendo singolarmente a un minuto di distanza l'uno dall'altro, e di procedere tutti con lo stesso passo costante di 5:00 al~km, in modo da non incontrarsi. Antonio, invece, alieno alle regole, è partito in un momento a caso tra il 30-esimo e il 31-esimo, ed ha corso ad un passo costante ma diverso dagli altri.

Sapendo che Antonio ha effettuato 6 sorpassi nei 10~km, cosa possiamo dire del suo passo?

\item  Dopo due mesi, con il diminuire delle restrizioni, i concorrenti hanno confermato le partenze singole distanziate di un minuto, ma hanno deciso che ognuno potesse scegliere se procedere con il passo costante di 4:50, 5:10 o 5:30 al~km. Supponendo che ognuno abbia scelto con uguale probabilità tra le tre opzioni, indipendentemente dalle scelte altrui, determinare la probabilità che il concorrente partito per 30-esimo abbia effettuato e subito lo stesso numero di sorpassi.

\item  All'ultimo allenamento, grazie all'ulteriore allentamento delle restrizioni, i componenti del gruppo hanno deciso di partire a coppie, sempre a distanza di un minuto dalla coppia precedente. In ogni coppia, uno correva al passo costante di 5:00 al~km e uno al passo costante di 5:20 al~km. Ancora una volta Antonio non ha rispettato la regola, partendo in un momento a caso tra il 30-esimo e il 31-esimo gruppo, e tenendo un passo costante diverso dagli altri. A fine allenamento, Antonio ha dichiarato: ``ho effettuato 1 sorpasso e ne ho subiti 2, ma se fossi partito 5 secondi dopo avrei effettuato 2 sorpassi e ne avrei subiti 2''.

Cosa possiamo dire del passo che ha tenuto Antonio?

\end{enumerate}

{\small Nota. Si intende che i sorpassi effettuati sulla linea di partenza/arrivo non si contano. Il passo di un corridore è di 5:27 al~km se il corridore impiega 5 minuti e 27 secondi per percorrere 1~km.}
}

\renewcommand{\Risposta}{
% SSSUP-BOZZA-AI: Risposta; autore=ChatGPT; data=2026-11-08; stato=DA-VERIFICARE
% SSSUP-PROVENIENZA: AI: elaborazione nuova dal testo
\begin{enumerate}
\item[(a)] Il passo è strettamente compreso tra $4{:}18$ e $4{:}30$ al km.
\item[(b)] La probabilità è $29/243$.
\item[(c)] Il passo è strettamente compreso tra $5{:}11{,}5$ e $5{:}12{,}5$ al km. Se lo si esprime in un numero intero di secondi al km, l'unica possibilità è $5{:}12$.
\end{enumerate}
}

\renewcommand{\Aiutino}{
% SSSUP-BOZZA-AI: Aiutino; autore=ChatGPT; data=2026-11-08; stato=DA-VERIFICARE
% SSSUP-PROVENIENZA: AI: elaborazione nuova dal testo
Con velocità costanti, due corridori si sorpassano durante il percorso esattamente quando l'ordine di arrivo è opposto all'ordine di partenza, escludendo le uguaglianze all'arrivo. Lavora sui tempi totali dei $10$ km. Nel punto (c), partire cinque secondi dopo può aumentare il numero di sorpassi solo passando dall'altra parte della partenza di una coppia.
}

\renewcommand{\SoluzioneUfficiale}{}

\renewcommand{\SoluzioniAlternative}{}

\renewcommand{\Osservazioni}{}

\renewcommand{\Commenti}{}

% Aggiunta redazionale: la soluzione preesistente non viene modificata.
\providecommand{\SoluzioneRedazionale}{}
\renewcommand{\SoluzioneRedazionale}{
% SSSUP-BOZZA-AI: SoluzioneRedazionale; autore=ChatGPT; data=2026-11-08; stato=DA-VERIFICARE
% SSSUP-PROVENIENZA: AI: elaborazione nuova dal testo
Misuriamo i tempi in minuti. Per due corridori che corrono a passo costante, un sorpasso nel tratto aperto fra partenza e arrivo equivale all'inversione stretta dell'ordine fra partenza e arrivo.

\paragraph{Punto (a).}
Prendiamo come tempo zero la partenza del trentesimo corridore e indichiamo con $t\in(0,1)$ il ritardo di Antonio. Se $T$ è il suo tempo sui $10$ km, poniamo $D=50-T$. Il corridore partito al tempo $-j$, con $j=0,1,\ldots,29$, viene superato esattamente quando
\[
t+T<50-j,\qquad\text{cioè}\qquad j<D-t.
\]
Avere esattamente sei sorpassi equivale a $5<D-t\leq6$. Poiché $0<t<1$, si ha $5<D<7$, e ogni valore interno è realizzabile. Quindi $43<T<45$ minuti: dividendo per $10$ si ottiene l'intervallo della risposta.

\paragraph{Punto (b).}
I tempi dei tre ritmi sono $48$ minuti $20$ secondi, $51$ minuti $40$ secondi e $55$ minuti. Fra ritmi consecutivi la differenza è $3$ minuti $20$ secondi, fra gli estremi $6$ minuti $40$ secondi.

Se il trentesimo è veloce, non subisce sorpassi. Per non effettuarne deve trovare il ritmo veloce nei tre predecessori immediati e un ritmo non lento nei tre ancora precedenti. La probabilità condizionata è $(1/3)^3(2/3)^3=8/729$. Se è lento, per simmetria, si ottiene la stessa probabilità, guardando i sei successori.

Se ha ritmo intermedio, supera i lenti fra i tre predecessori immediati ed è superato dai veloci fra i tre successori immediati. I due conteggi sono indipendenti e hanno distribuzione binomiale con parametri $3,1/3$. La probabilità che siano uguali è
\[
\sum_{j=0}^{3}\binom3j^2
\left(\frac13\right)^{2j}\left(\frac23\right)^{6-2j}
=\frac{64+144+36+1}{729}=\frac{245}{729}.
\]
Mediando sui tre ritmi si trova $(8+245+8)/(3\cdot729)=29/243$.

\paragraph{Punto (c).}
I due tempi dei componenti delle coppie sono $50$ e $160/3$ minuti. Antonio deve avere $50<T<160/3$, altrimenti non potrebbe sia effettuare sia subire sorpassi. Poniamo
\[
D_f=T-50,\qquad D_s=160/3-T,\qquad D_f+D_s=10/3,
\]
e ancora $t\in(0,1)$ come ritardo rispetto alla trentesima coppia. Un sorpasso effettuato significa $0<D_s-t\leq1$; due sorpassi subiti significano $2<t+D_f\leq3$.

Ritardare di $\delta=1/12$ di minuto, restando tra le medesime due coppie, non può aumentare i sorpassi effettuati. Deve quindi essere $11/12<t<1$: Antonio parte, nella situazione ipotetica, dopo la coppia successiva. Per continuare a subire esattamente due sorpassi, perduto quello del veloce di tale coppia, deve aggiungersi il sorpasso del veloce della terza coppia successiva. Ponendo $q=t+D_f$, la condizione è
\[
\frac{35}{12}<q\leq3.
\]
Queste condizioni garantiscono anche i due sorpassi effettuati: prima il lento della trentesima coppia, e dopo il ritardo i lenti della trentesima e trentunesima, senza raggiungere il precedente.

Eliminando $t$ da $11/12<t<1$ e $35/12<t+D_f\leq3$ si ottiene
\[
\frac{23}{12}<D_f<\frac{25}{12},\qquad
51\text{ min }55\text{ s}<T<52\text{ min }5\text{ s}.
\]
Ogni valore interno ammette una scelta di $t$ che rispetta le condizioni. Dividendo per $10$ si conclude.
}
